\begin{finding}[\fid{F28}: an isolated hook excludes all Young-stabilizer monomial models]\label{f:recthooks}
Let $M$ be a characteristic-zero $S_N$-module. If an internal hook
degree $1\le k\le N-2$ occurs in $M$, while its two neighboring
hook degrees do not, then $M$ has no scalar-monomial basis all of
whose label stabilizers are conjugate to Young subgroups.

Consequently, for \emph{every} rectangle $\rho=(a^b)$ with $a,b\ge2$,
the actual module $[\rho]\otimes[\rho]$ has no such basis, even when
arbitrary nonzero scalar phases are allowed. Thus the Young-stabilizer
version of \fid{F23} cannot describe the unprojected rectangular
Satake multiplicity module, including large-parameter rectangles.
This does not exclude non-Young stabilizers or other positive rules.
\end{finding}
\begin{proof}
For a Young subgroup $H=\prod_{i=1}^r S_{a_i}$ with nonempty blocks,
every one-dimensional character is a product of trivial and sign
characters. Set $s=\sum_{\text{sign blocks }i}(a_i-1)$; singleton
blocks contribute zero under either designation. If $U_N$ is the
standard sum-zero representation, then
\[
 \Res_H U_N=\bigoplus_i U_{a_i}\oplus k^{r-1},\qquad U_1=0.
\]
The classical identification $\wedge^jU_a=[(a-j,1^j)]$ follows here
from the cited Pieri rule: the oriented $j$-subset module
$\wedge^j k^a=\operatorname{Ind}_{S_j\times S_{a-j}}^{S_a}
(\mathrm{sgn}\boxtimes1)$ is the sum of the two adjacent hooks for $1\le j\le a-1$,
and $k^a=1\oplus U_a$ identifies their exterior powers inductively.
In the exterior power of the displayed restriction, trivial on block
$i$ forces degree zero, and sign forces degree $a_i-1$. The remaining
degree is chosen from the trivial $(r-1)$-space. Frobenius reciprocity
therefore gives the exact formula
\[
 [\operatorname{Ind}_H^{S_N}\chi:(N-j,1^j)]
       =\binom{r-1}{j-s}.
\]
Its positive hook support is the consecutive interval
$s,\ldots,s+r-1$, contained in $0,\ldots,N-1$ because
$0\le s\le N-r$.

A scalar-monomial basis decomposes into actual orbit modules
$\operatorname{Ind}_{H_O}^{S_N}\chi_O$, by the orbit-map proof in
\fid{F23}. Some summand must contain the given internal hook. If its
Young subgroup has $r\ge2$ blocks, the interval just computed includes
at least one neighboring hook, contradicting its absence from the
effective direct sum. If $r=1$, the orbit is the trivial or sign
representation of the whole group, whose hook degree is $0$ or
$N-1$, not internal. This proves the first assertion. There is no
cancellation in an effective representation decomposition.

The exact source input for rectangles is Ikenmeyer--Panova
\cite[Proposition~6.4, pp.~15--16]{IP}. With $m=\min(a,b)$ and $N=ab$,
it states
\begin{equation}\label{eq:recthookproduct}
 \sum_{j=0}^{N-1}g((N-j,1^j),\rho,\rho)t^j
       =\prod_{i=2}^{m}(1+t^{2i-1}).
\end{equation}
Transposition of both rectangular Specht factors justifies the
symmetric shorter-side convention. For $N\ge6$, degree three is
internal. The product has coefficient one at degree three and zero
at degrees two and four. Apply the isolated-hook assertion.
The only smaller rectangle is $(2,2)$ of size four. Its square is
$[(4)]\oplus[(2,2)]\oplus[(1^4)]$ by the already checked character
calculation, or Manivel's Theorem~1. Its hook degrees are $0,3$.
Every Young orbit with at least two blocks would force consecutive
hooks; all remaining whole-group orbits are one-dimensional and
cannot supply $[(2,2)]$. This proves the edge case as well.
\end{proof}

The actual multiplicity identification is explicit: for $\dim U=\dim V=b$,
Schur--Weyl duality gives
\[
 \Hom_{\GL(U)\times\GL(V)}
 (\det^aU\boxtimes\det^aV,(U\otimes V)^{\otimes ab})
 \simeq[(a^b)]\otimes[(a^b)]
\]
with its diagonal symmetric-group action. Its primitive Schur
projectors give precisely the rectangular Kronecker coefficients.
The obstruction thus concerns the requested actual family, not an
unrelated test representation. A virtual alternating induction formula
can cancel the forbidden hooks, but would not be the positive orbit
decomposition considered here.

The scope is sharp already at size four. The permutation module on
the three pairings of four letters is $1\oplus[(2,2)]$, so
$[(2,2)]^{\otimes2}$ is that module plus a sign line: it does have
a signed-permutation model. A pairing stabilizer is
$S_2\wr S_2$, of order eight, whereas a Young subgroup of $S_4$ has
order $24,6,4,2$ or $1$. The stabilizer is not Young.

\paragraph{A working ordinary monomial model in the known height-two family.}
Manivel \cite[Theorem~1, p.~1; \S\S2--3, pp.~2--4]{Manivel} proves
that $[(d,d)]^{\otimes2}$ is multiplicity free, with constituents the
partitions of $2d$ having at most four even parts, or exactly four odd
parts. In fact, for dimensions $(\dim U,\dim V,\dim W)=(2,2,4)$,
his proof identifies the highest-weight invariant algebra
\[
 \C[U\otimes V\otimes W]^{SL(U)\times SL(V)\times N_W}
 =\C[f_1,f_2,f_3,f_4]
\]
in the source's dual-polynomial convention. The four $W$-partition
labels are $(1^4)$, $(2)$, $(2,2)$ and $(2,2,2)$, with degrees $4,2,4,6$.
These are known source results, not consequences newly obtained from
purity. A monomial with label $\nu$, padded to four parts, necessarily
has exponents
\[
 a_1=\nu_4,\quad a_2=(\nu_1-\nu_2)/2,\quad
 a_3=(\nu_2-\nu_3)/2,\quad a_4=(\nu_3-\nu_4)/2.
\]
They are nonnegative by the partition inequalities, and integral
exactly when the padded parts have common parity. The two determinant
exponents are $d=2a_1+a_2+2a_3+3a_4=|\nu|/2$.
Thus this is an actual known monomial basis with a verified
multiplicity comparison, in a different representation-theoretic model.

Equivalently, for $\ell(\nu)\le4$, let
\[
 P_\nu=\{((\nu_1-\nu_4)/2,(\nu_2-\nu_4)/2,(\nu_3-\nu_4)/2)\}.
\]
Its coordinates are $(a_2+a_3+a_4,a_3+a_4,a_4)$, an integral
unimodular change from the last three exponents. The exact restricted
Ehrhart formula is
\[
 g(N\nu,(Nd,Nd),(Nd,Nd))=\#(NP_\nu\cap\Z^3).
\]
Use the empty polytope for length greater than four. The point has
denominator at most two and a polynomial-bit description and
separation test. It is the source's actual exponent rule, not a
point assigned after consulting a zero/nonzero oracle. Its torus
monomial basis, for $d\ge2$, is not a scalar-monomial $S_{2d}$ Young-orbit basis;
the latter is exactly what \fid{F28} excludes.


